% author: Matheus T. dos Santos
% co-authors:
%   - Claude (claude-opus-4-8), 2026-08-13 -- formatacao mecanica autorizada por Matheus
%     (suspensao pontual da Regra 3): italico em estrangeirismos (todas as ocorrencias) no
%     corpo do resumo. Siglas nao convertidas para \ac{} (o resumo precede a Lista de Siglas).
%     Nenhuma palavra do texto foi alterada.
\begin{resumo}

Em controle digital de malha fechada, não cumprir o requisito de tempo do período de amostragem pode desestabilizar o sistema por completo, causando acidentes graves. A natureza do controle digital exige estados compartilhados: seja na aquisição de dados, seja na alteração de parâmetros feita pelo usuário, seja na atuação do sinal de controle. Os estados compartilhados acessados por diversos contextos, onde pelo menos um deles realiza uma escrita e pelo menos um acesso é não sincronizado, levam a um \textit{data race}. O estado da arte hoje na implementação de controle digital, em sistemas embarcados, envolve C+MISRA+Sanitizers para evitar \textit{data races}. Entretanto este método depende da disciplina do desenvolvedor e de detectar \textit{data races} exercitados em tempo de execução. Com a linguagem Rust, essas garantias são movidas do desenvolvedor para o compilador e os \textit{data races} não são mais detectados em tempo de execução, mas sim eliminados em tempo de compilação, através do sistema de tipos do Rust (\textit{Send}/\textit{Sync} e \textit{borrow checker}). Todavia, o Rust exige uma sincronização forçada a qual possui um custo. Também é desconhecido quais resquícios de \textit{unsafe} ainda permanecem no Rust. Deste modo, qual a fronteira entre \textit{safe} e \textit{unsafe} nos padrões de \textit{data race}, no domínio de algoritmos embarcados de controle, e qual o custo das garantias do Rust? Neste trabalho, é proposta uma taxonomia para classificar os padrões de \textit{data race} presentes nos algoritmos de controle. Após isso, a fronteira entre \textit{safe} e \textit{unsafe} nestes padrões é caracterizada. Além disso, é catalogado o espaço de design das garantias que tornam o \textit{data race} inexprimível. Também é construído o veículo de experimentos: Aule, uma biblioteca de controle construída para instanciar os algoritmos de controle utilizados nos experimentos. Atualmente, Aule contém a base principal para a construção dos experimentos em uma planta física, com os blocos de sincronização em desenvolvimento. Também está presente neste trabalho um panorama comparativo do ecossistema \textit{open-source} do Rust para bibliotecas de controle. Por fim, é mostrado neste trabalho um caso demonstrativo de recusa do compilador Rust \textit{safe}, quando o código possui \textit{data race}.

\palavrasChave{\textit{Data race}; Segurança de Memória; Sistemas de Controle; Rust; Sistemas Embarcados}

\end{resumo}
